\chapter*{Résumé}
\addcontentsline{toc}{chapter}{Résumé}

Cette thèse développe des architectures neuronales et des méthodes d'optimisation efficaces basées sur l'orthogonalité. Elle part des réseaux de neurones 1-Lipschitziens, où les couches orthogonales offrent un contrôle certifiable de la robustesse, et résout les problèmes de flexibilité, de passage à l'échelle et de qualité d'implémentation. Elle étend ensuite la perspective de l'orthogonalité à l'optimisation, en étudiant des méthodes de type Muon. Les contributions comprennent de nouvelles couches convolutives orthogonales, une bibliothèque open-source de composants 1-Lipschitziens, une variante accélérée de Muon et une étude des méthodes d'optimisation orthogonale pour les convolutions.
